\par\medskip\pdfbookmark[1]{On V.}{nexus-note-V}\centerline{On V.}\smallskip
The remark appended to formula III indicates how the author arrived at the number $k$ of forms in the collection $C$. Geometrical considerations (compare I, art. 5, and II, art. 4) give the following limit for the quotient
\[
\frac{(1)+(2)+(3)+\ldots+(m)}m
\]
as $m$ grows without bound, that is, the average number of representations of an unspecified positive integer,
\[
=k\frac\pi{\surd D}
\]
A second expression for this same limit can be derived in several ways from the nature of the number $(m)=\epsilon f(m)$ of representations of the number $m$, determined above. The surviving fragments (I, arts. 7, 8; III and IV) suggest that the author's first approach was as follows.

Let $\theta(m)$ be any function of the positive integer $m$, and $p$ any prime. From $\theta(m)$ one can always derive a new function $\theta'(m)$ whose value is independent of whether, and how often, $p$ occurs as a factor in $m$, and which, whenever $p$ does not divide $m$, agrees with $\theta(m)$. Such a function is obtained by putting $\theta'(m)=\theta\left(\frac m{p^\pi}\right)$, where $p^\pi$ denotes the highest power dividing $m$ with base $p$; one may say that the function $\theta'(m)$ arises from $\theta(m)$ by elimination of the prime $p$. If in this way from $f(m)$ one forms a new function $f'(m)$ by eliminating the prime $2$, then forms $f''(m)$ from it by eliminating $3$, and so on, each successive function will have a more regular course than its predecessors. If the primes are eliminated one after another in increasing order, such a function
\SourceLeaf{305}
$\theta(m)$ will, for infinitely many values of $m$, have the value $f(1)=1$, in particular for all values of $m$ smaller than the last prime eliminated. Continuing this process indefinitely brings one ever closer to the function $f^\infty(m)$, which for all values of $m$ has the value $1$, and whose average value therefore also $=1$. If it is possible to express the average value of any function $\theta(m)$ in terms of that of its successor $\theta'(m)$, it will also be possible to find the average value of $f(m)$ through an infinite chain of operations.

If $p$ is the prime whose elimination produces $\theta'(m)$ from $\theta(m)$, then $\theta(m)=\theta'(m)f(p^\pi)$, where $p^\pi$ again denotes the highest power dividing $m$ with base $p$. When $p$ does not divide $2D$, this readily gives
\[
\theta'(m)=\theta(mp)-\left(\frac{-D}p\right)\theta(m)
\]
Writing for brevity
\[
\begin{aligned}
\Theta(m)&=\theta(1)+\theta(2)+\ldots+\theta(m)\\
\Theta'(m)&=\theta'(1)+\theta'(2)+\ldots+\theta(m)
\end{aligned}
\]
one obtains
\[
\Theta'(mp)=\Theta(mp)-\left(\frac{-D}p\right)\Theta(m)
\]
and hence, using $\omega,\omega'$ for the respective average values of the functions $\theta(m),\theta'(m)$,
\[
\omega=\frac{\omega'}{1-\left(\frac{-D}p\right)\frac1p}
\]

If, however, the prime $p$ divides $2D$, the functions $\theta(m)$ and $\theta'(m)$ no longer have so simple a relation in general; nevertheless, a similar argument readily gives $\omega=\omega'$ in this case. The relation between $\omega$ and $\omega'$ can also be derived in both cases as follows. Put
\[
\vartheta(m)=\Sigma\theta(\mu)
\]
where the sum extends over all numbers $\mu$ that are not divisible by $p$ and are no greater than $m$. Denote by $m',m'',m'''\ldots$, respectively, the greatest integers contained in $\frac mp,\frac{m'}p,\frac{m''}p\ldots$. Then
\[
\begin{aligned}
\Theta(m)&=\vartheta(m)+\vartheta(m')f(p)+\vartheta(m'')f(pp)+\vartheta(m''')f(p^3)+\ldots\\
\Theta'(m)&=\vartheta(m)+\vartheta(m')+\vartheta(m'')+\vartheta(m''')+\ldots
\end{aligned}
\]
and this gives
\[
\frac\omega{\omega'}=(1-\tfrac1p)\{1+\frac{f(p)}p+\frac{f(pp)}{pp}+\frac{f(p^3)}{p^8}+\ldots\}
\]
in agreement with the result found above (compare the note on III and IV).

The average value of the function $f(m)$ is therefore equal to the infinite product
\[
\Pi\frac1{1-\left(\frac{-D}p\right)\frac1p}
\]
\SourceLeaf{306}
in which $p$ must run through all primes not dividing $2D$, and hence
\[
k=\frac{\epsilon\surd D}\pi\Pi\frac1{1-\left(\frac{-D}p\right)\frac1p}
\]

There remains, however, a concern about the rigor of this deduction, concerning the method of determining the average value of $f(m)$ by successively eliminating all primes. For it is evident that after elimination of the first $n$ primes, the resulting function $f^{(n)}(m)$ agrees in value with the function $f^\infty(m)=1$ as long as $m$ remains smaller than the last prime eliminated. Thus, by choosing $n$ large enough, this agreement can be extended up to any prescribed size of the number $m$. But this alone by no means proves that, as $n$ grows without bound, the average value of $f^{(n)}(m)$ approaches without limit the average value of $f^\infty(m)$, that is, the value $1$. How the author intended to fill this gap cannot be determined with certainty from the surviving papers. Yet the formula already given above (in the note on I),
\[
1+4\{A-\frac A3+\frac A5-\frac A7+\frac A9-\frac A{11}\ldots\}
\]
for the number of pairs of numbers whose sum of squares does not exceed $A$, suggests that the author bypassed the infinite product and directly found, for the average value of $f(m)$, the infinite series
\[
\Sigma\left(\frac{-D}n\right)\frac1n
\]
in which $n$ must run in increasing order through all positive integers relatively prime to $2D$. The simplest way of indicating this transition seems to be the following.

\clearpage
Let $\mu$ be the greatest divisor of a number $m$ among those relatively prime to $2D$, and put $\theta(m)=f(\mu)$. Then $\theta(m)$ is the function obtained by eliminating all primes dividing $2D$ from $f(m)$; according to the preceding discussion, its average value agrees with that of $f(m)$. Since $\theta(m)=\Sigma\left(\frac{-D}n\right)$, where $n$ runs through all divisors of $\mu$, that is, all divisors of $m$ relatively prime to $2D$, one obtains the formula analogous to the one above,
\[
\Theta(m)=\theta(1)+\theta(2)+\ldots+\theta(m)=\Sigma\left(\frac{-D}n\right)\frac mn
\]
where, in the sum on the right, the letter $n$ runs through all numbers relatively prime to $2D$, and only the greatest integer contained in the quotient $\frac mn$ is to be retained. Arrange the terms of this series so that the numbers $n$ follow in increasing order. The factor $\frac mn$ then continually decreases, or at least never increases, and the series terminates as soon as $n>m$. Furthermore, the fundamental theorem in the theory of quadratic residues and its generalization imply that the sum of every $\varphi(4D)$ consecutive values of the factor $\left(\frac{-D}n\right)$ vanishes. It follows that the sum of any number of consecutive values
\SourceLeaf{307}
of this factor never exceeds, in absolute value, the finite bound depending only on the determinant $D$, namely $\Delta=\varphi(2D)$. Combining these two observations, one readily finds that the sum of all terms following the term $\left(\frac{-D}n\right)\frac mn$ is smaller in absolute value than $\Delta\frac mn$, and consequently that the quotient $\Theta(m):m$, as $m$ grows without bound, has as its limit the convergent infinite series
\[
\Sigma\left(\frac{-D}n\right)\frac1n
\]
ordered in the way specified. Having thus found the common average value of the functions $\theta(m)$ and $f(m)$, one immediately obtains
\[
k=\frac{\epsilon\surd D}\pi\Sigma\left(\frac{-D}n\right)\frac1n
\]

It should also be observed that Articles 6 and 8 of paper II point to an approach to the problem that is simpler in some respects and also easy to carry out, in which only representations of odd numbers, or even only of numbers relatively prime to $2D$, are considered.
\BookRule{35mm}
